% from ega2-indexes-and-contents-fr.tex
\oldpage[II]{210}
\markboth{A. GROTHENDIECK --- Chap. II}{A. GROTHENDIECK --- Chap. II}
\egaTerme{Fibre géométrique rationnelle d'un fibré projectif sur une extension $K$ de $k(y)$}{\hyperref[II.4.2.6-fr]{4.2.6}}
\egaTerme{Fibré projectif défini par un $\mathcal{O}_Y$-Module}{\hyperref[II.4.1.1-fr]{4.1.1}}
\egaTerme{Fibré vectoriel défini par un $\mathcal{O}_S$-Module}{\hyperref[II.1.7.8-fr]{1.7.8}}
\egaTerme{Fini (morphisme)}{\hyperref[II.6.1.1-fr]{6.1.1}}
\egaTerme{Fini (préschéma) au-dessus d'un préschéma}{\hyperref[II.6.1.1-fr]{6.1.1}}
\egaTerme{Finie ($\mathcal{O}_S$-Algèbre quasi-cohérente)}{\hyperref[II.6.1.2-fr]{6.1.2}}
\egaTerme{Fonctions symétriques élémentaires d'un endomorphisme d'un module localement libre}{\hyperref[II.6.4.1-fr]{6.4.1}}
\egaTerme{Fonctions symétriques élémentaires d'un endomorphisme d'un module de type fini sur un anneau intègre}{\hyperref[II.6.4.2-fr]{6.4.2}}
\egaTerme{Fonctions symétriques élémentaires d'un endomorphisme d'un module de type fini sur un anneau noethérien réduit}{\hyperref[II.6.4.7-fr]{6.4.7}}
\egaTerme{Fonctions symétriques élémentaires d'un endomorphisme d'un $\mathcal{A}$-Module localement libre de rang constant}{\hyperref[II.6.4.8-fr]{6.4.8}}
\egaTerme{Fonctions symétriques élémentaires d'un endomorphisme d'un $\mathcal{O}_X$-Module de type fini sur un préschéma $X$ localement intègre}{\hyperref[II.6.4.9-fr]{6.4.9}}
\egaTerme{Fonctions symétriques élémentaires d'un endomorphisme d'un $\mathcal{O}_X$-Module de type fini sur un préschéma $X$ localement noethérien réduit}{\hyperref[II.6.4.10-fr]{6.4.10}}
\egaTerme{Fonctions symétriques élémentaires d'une section d'une $\mathcal{A}$-Algèbre}{\hyperref[II.6.5.1-fr]{6.5.1}}

\egaTerme{Homogénisé d'un sous-$\mathcal{S}$-Module d'un $\mathcal{S}$-Module gradué}{\hyperref[II.8.13.2-fr]{8.13.2}}
\egaTerme{Homomorphisme d'anneaux gradués}{\hyperref[II.2.1.2-fr]{2.1.2}}
\egaTerme{Homomorphisme de degré $k$ de modules gradués}{\hyperref[II.2.1.2-fr]{2.1.2}}

\egaTerme{Idéal de $S_+$, idéal premier gradué de $S_+$ ($S$ anneau gradué)}{\hyperref[II.2.1.10-fr]{2.1.10}}

\egaTerme{Lemme de Chow}{\hyperref[II.5.6.1-fr]{5.6.1}}
\egaTerme{Lieu à l'infini d'un cône projectif}{\hyperref[II.8.3.3-fr]{8.3.3}}

\egaTerme{Morphisme canonique $G(\varepsilon)\to\Proj(S)$, $S=\bigoplus_{n\geq0}\Gamma(X,\mathcal{L}^{\otimes n})$}{\hyperref[II.4.5.1-fr]{4.5.1}}
\egaTerme{Morphisme canonique $X\to\Spec(\Gamma(X,\mathcal{O}_X))$}{\hyperref[II.5.1.1-fr]{5.1.1}}
\egaTerme{Morphisme de Segre}{\hyperref[II.4.3.1-fr]{4.3.1}}

\egaTerme{Nilradical de $S_+$ ($S$ anneau gradué)}{\hyperref[II.2.1.10-fr]{2.1.10}}
\egaTerme{Nilradical de $\mathcal{S}_+$ ($\mathcal{S}$ $\mathcal{O}_Y$-Algèbre graduée quasi-cohérente)}{\hyperref[II.3.1.12-fr]{3.1.12}}
\egaTerme{Norme d'une section d'une $\mathcal{A}$-Algèbre}{\hyperref[II.6.5.1-fr]{6.5.1}}
\egaTerme{Norme d'un $\mathcal{B}$-Module inversible}{\hyperref[II.6.5.2-fr]{6.5.2}}
\egaTerme{Norme d'une section d'un $\mathcal{B}$-Module inversible}{\hyperref[II.6.5.3-fr]{6.5.3}}
\egaTerme{Norme d'un $\mathcal{O}_{X'}$-Module inversible}{\hyperref[II.6.5.5-fr]{6.5.5}}

\egaTerme{Périodique (Algèbre graduée)}{\hyperref[II.8.14.12-fr]{8.14.12}}
\egaTerme{Polynôme caractéristique d'un endomorphisme d'un module de type fini sur un anneau intègre}{\hyperref[II.6.4.2-fr]{6.4.2}}
\egaTerme{Polynôme caractéristique d'un endomorphisme d'un module de type fini sur un anneau noethérien réduit}{\hyperref[II.6.4.7-fr]{6.4.7}}
\egaTerme{Préschéma des sommets d'un cône affine (projectif)}{\hyperref[II.8.3.3-fr]{8.3.3}}
\egaTerme{Présentation finie (module gradué ayant une)}{\hyperref[II.2.1.1-fr]{2.1.1}}
\egaTerme{Présentation finie ($\mathcal{S}$-Module gradué ayant une)}{\hyperref[II.3.1.1-fr]{3.1.1}}
\egaTerme{Produit tensoriel gradué de deux modules gradués}{\hyperref[II.2.1.2-fr]{2.1.2}}
\egaTerme{Projectif (morphisme)}{\hyperref[II.5.5.2-fr]{5.5.2}}
\egaTerme{Projectif (préschéma) au-dessus d'un préschéma}{\hyperref[II.5.5.2-fr]{5.5.2}}
\egaTerme{Propre (morphisme)}{\hyperref[II.5.4.1-fr]{5.4.1}}
\egaTerme{Propre (partie)}{\hyperref[II.5.4.10-fr]{5.4.10}}
\egaTerme{Propre (préschéma) au-dessus d'un préschéma}{\hyperref[II.5.4.1-fr]{5.4.1}}

\egaTerme{Quasi-affine (morphisme)}{\hyperref[II.5.1.1-fr]{5.1.1}}
\egaTerme{Quasi-affine (préschéma) au-dessus d'un préschéma}{\hyperref[II.5.1.1-fr]{5.1.1}}
\egaTerme{Quasi-affine (schéma)}{\hyperref[II.5.1.1-fr]{5.1.1}}

